\section{Reception: Relics, Doctor, and the Contested Inheritance}

\subsection{Earliest memory}

The cult began with the biography. Possidius, writing at Calama within a few
years of the death, already proposed his friend for imitation and preserved
the \work{Indiculus} so that the works could be sought ``from the library of
the church of Hippo, where the more perfect copies can probably be found''
(\work{Vita} 18). Within about a year of the death, as John Paul II noted
for the sixteenth centenary of the conversion, Pope Celestine I was already
calling Augustine ``one of the best teachers'' of the Church
(\work{Augustinum Hipponensem}, 1986). The feast stands on 28 August, the
recorded day of death.

\subsection{The relics: Hippo, Sardinia, Pavia}

Augustine was buried at Hippo in 430. The next secure sighting of his body
is nearly three centuries later and five hundred miles away. Paul the
Deacon, the Lombard historian writing in the later eighth century, records
under the reign of King Liutprand (712--744): ``Liutprand also, hearing that
the Saracens had laid waste Sardinia and were even defiling those places
where the bones of the bishop St.~Augustine had been formerly carried on
account of the devastation of the barbarians and had been honorably buried,
sent and gave a great price and took them and carried them over to the city
of Ticinum''---Pavia---``and there buried them with the honor due to so
great a father'' (\work{History of the Langobards} 6.48). The same sentence
is the earliest witness located by this research for the first translation
as well: the bones had been ``formerly carried'' to Sardinia because of
``the barbarians,'' a removal later tradition attributes to the African
bishops exiled by the Vandal kings around 500, though Paul names no one and
gives no date. The medieval legendaries embroidered the second translation
---the \work{Golden Legend} has ``one Liprand, a devout king of the
Lombards,'' pay ``great good'' for the body, meet it at Genoa, and found a
church at every overnight stop until it was ``laid honourably in the church
of S.~Peter which is called Cieldore, or heaven of gold.'' There, in San
Pietro in Ciel d'Oro at Pavia, the shrine remains: Benedict XVI, preaching
in that city in 2007, recalled that ``the King of the Longobards acquired
Augustine's remains for the City of Pavia so that today they belong to this
City in a special way'' (Homily at Pavia, 22 April 2007). Monica's remains,
by a tradition reported in the reference literature, were brought from
Ostia to Rome in 1430 under Martin V and rest in the church of
Sant'Agostino.

\witnesslimit{Between the burial of 430 and Paul the Deacon's notice, no
contemporary record of the relics' movements is preserved: the
Sardinia translation is inference from an eighth-century sentence plus
later tradition, and the identification of the Pavia bones with the body
buried at Hippo rests on that same chain. The medieval and early-modern
recognitions of the Pavia relics, and their modern arrangement in the Arca
of San Pietro in Ciel d'Oro, are matters of record for their own dates
only; this study asserts the continuous chain as received tradition, not
as verified identification, while noting that the shrine's standing in
current official usage (the papal visit of 2007) is itself a fact of
reception.}

\subsection{Doctor of Grace}

Augustine's doctrinal authority in the Latin church was quasi-canonical
from the fifth century onward. The medieval church formalized what custom
held: Boniface VIII's decretal of 1298 ordered the feasts of the four great
Western teachers---Gregory the Great, Ambrose, Augustine, and Jerome---kept
throughout the Church, the act from which the title ``Doctor of the
Church'' descends (so the 1912 Catholic Encyclopedia, ``Doctors of the
Church''). The tradition's particular name for him states his
specialization: ``he is not only the Doctor of Grace, he is also the Doctor
of the Church,'' as the older Catholic Encyclopedia puts it in summarizing
his double legacy. The modern papacy has returned to him at length: John
Paul II with the apostolic letter \work{Augustinum Hipponensem} (1986) for
the sixteenth centenary of the conversion, calling him ``the common father
of Christian Europe,'' and Benedict XVI with the Pavia pilgrimage of 2007
and five general audiences of January and February 2008 devoted to the one
Father. The Catechism reaches for him
on its opening pages, closing its paragraphs on the desire for God with the
prayer of the restless heart (CCC 30).

\subsection{Medieval and Reformation appropriations}

For the Latin Middle Ages Augustine was simply ``the theologian'': the
monastic centuries read God through his Psalms commentaries; the canons
regular and, from the thirteenth century, the Order of Hermits of
St.~Augustine lived under a rule bearing his name (its textual history---a
short rule with genuine Augustinian material and a complicated
transmission---is a specialist question this study does not adjudicate);
the scholastics cited him as the standing authority whom Aristotle had to
be reconciled with, not against. The Reformation then fought over the
inheritance. Luther was an Augustinian friar; he and Calvin claimed the
anti-Pelagian Augustine of grace and predestination against what they took
to be a new Pelagianism, while the Council of Trent answered with the same
doctor's teaching on baptism, justification, and the Church. The old
Encyclopedia's summary of the Protestant historians still serves:
``Luther and Calvin were content to treat Augustine with a little less
irreverence than they did the other Fathers,'' while Harnack could ask
``Where, in the history of the West, is there to be found a man who, in
point of influence, can be compared with him?'' (Catholic Encyclopedia,
``Teaching of St.~Augustine of Hippo,'' quoting Harnack). Jansenism in the
seventeenth century claimed his hardest anti-Pelagian pages against the
Church that had canonized him, and was condemned without any censure
falling on Augustine himself---the ``Jansenistic abuse of his works,'' as
the same survey calls it, alarming some without changing the general
verdict: the Church has consistently received the Doctor of Grace's
doctrine while declining to bind itself to every polemical inference of his
last years.

\subsection{The modern audit}

Modern critical scholarship rebuilt the man inside his province. The
standard account in English remains Peter Brown's \work{Augustine of Hippo:
A Biography} (1967), reissued in 2000 with ``a substantial Epilogue''
reconsidering the portrait in the light of ``the remarkable discovery of a
considerable number of letters and sermons by Augustine''---the Divjak
letters and Dolbeau sermons, which show the aged bishop at pastoral work in
unguarded detail; the publisher's description records the book's standing
as ``the standard account of Saint Augustine's life and teaching''
(University of California Press catalogue; the book itself is cited here
only at this record level). The larger movement of scholarship since has
run along lines this study has used at reported level: the Cassiciacum
dialogues weighed against the \work{Confessions}; the African, Punic, and
Berber context of Thagaste and Hippo taken seriously; the Donatists read,
so far as the sources allow, as an African church with its own theology of
martyrdom rather than a mere sect; and the late anti-Pelagian Augustine
treated as a problem the tradition itself has to manage. None of it has
displaced the two ancient witnesses; it has taught readers to hear how much
of what they know is Augustine's own voice, and to weigh it accordingly.

\subsection{Legends audited}

Three famous traditions require explicit status. First, the child by the
sea: Augustine, walking the shore while pondering the Trinity, meets a
child ladling the sea into a hole in the sand with a spoon and is told that
the sea will sooner fit the hole than the Trinity the human mind. The story
is absent from Augustine, Possidius, and the medieval legendary's own
sources: Caxton, printing his English \work{Golden Legend} in 1483, added
it from iconography with a bookseller's candor---``I have seen [it] painted
on an altar of S.~Austin at the black friars at Antwerp, howbeit I find it
not in the legend, mine exemplar, neither in English, French, ne in
Latin''---and thirteenth-century exemplum collections are reported as its
earliest literary carriers (a lead this research could not inspect). It is
late legend: theologically apt, historically unattested. Second, the
\work{Te Deum}: the \work{Golden Legend} reports, ``as it is read,'' that at
the baptism Ambrose and Augustine improvised the hymn in alternating
verses. The hymnological survey in the 1912 Catholic Encyclopedia traces
that tradition no earlier than the end of the eighth century and weighs the
competing ancient attributions---modern scholarship's candidate being
Niceta of Remesiana; the legend, not any contemporary witness, is why older
books call the \work{Te Deum} the ``Ambrosian hymn.'' Third, ``make me
chaste, but not
yet'': genuine, but a quotation of the penitent, not a quip of the sinner
---in context it is the mature bishop's example, told against himself, of
the divided will that only grace resolved (\work{Confessions} 8.7.17). The
pear tree, for its part, is no legend but the subject's own confession; its
status is that of everything else in that book---the past as grace
re-read it.
